\styledchapter{Collusion}

We briefly discuss collusion in the context of strong cartels, which allow for transfer payments between members of the bidding ring.

\section{Strong Cartels}

We consider an independent private values model with symmetric values $v_i \sim f(v_i)$ and $v_i \in [0,1]$ and all $N$ members participating in the bidding ring. Suppose the seller uses a first-price auction with reserve price $\rho$.\boxednote{We make no restriction on how $\rho$ is determined.}

The best the bidding ring can do win the good when $\max \left\{v_1,\ldots,v_N\right\} \ge \rho$, in which case they pay the seller $\rho$.
The bidders do this by running a ``knockout'' auction among themselves.

\begin{definition}[Knockout auction]
	In a knockout auction, the $N$ bidders in the bidding ring submit a sealed bid $b_i \in [0,\infty)$. The highest bid wins and the corresponding bidder earns the right to bid $\rho$ in the real auction, where all others must bid less than $\rho$ in the real auction. The high bidder keeps the good, and pays each of the other bidders the side payment $\frac{b-\rho}{N-1}$.
\end{definition}

The total payment made by a winner of the knockout auction (and the main auction) is $\rho + (N-1) \frac{b-\rho}{N-1} = b$. We see that this format is equivalent to a first price auction, except that losers get a side payment.

Suppose $\tilde{b}(\cdot )$ is a symmetric equilibrium bidding function for the knockout auction. Assume that $\tilde{b}$ is strictly increasing, with $\tilde{b}(v) = 0$ for $v < \rho$ and $\tilde{b}(\rho) = \rho$. We want to determine what $\tilde{b}$ is above $\rho$. Let $u(r,v)$ be a bidder's utility when her value is $v$ and she bids $\tilde{b}(r)$ and all others bid according to $\tilde{b}(\cdot )$.
We write down
\begin{align}
	u(r,v)
	&= F^{N-1}(r)\left( v - \tilde{b}(r) \right) \\
	&\quad + \int_{r}^{1}
		\underbrace{\left(\frac{\tilde{b}(s) - \rho}{N-1}\right)}_{\text{side payment}}
		\overbrace{\left( N-1 \right)F^{N-2}(s)f(s)}^{\mathclap{\text{density of highest bid among others}}}\,ds
.\end{align}

Maximizing as we have done, for $v \ge \rho$,
\begin{align*}
	\frac{du}{dr}
	= (N-1) F^{N-2}(r) f(r) \left( v- \tilde{b}(r) \right) &- F^{N-1}(r) \tilde{b}'(r) - \left( \tilde{b}(r) - \rho \right)F^{N-1}(r) f(r)
.\end{align*}
The equilibrium condition $\left.\frac{du}{dr}\right|_{r= v} = 0$ yields a linear ODE for $\tilde{b}$.
Multiplying by $F(v)$ and collecting terms gives \[
	G(v) = F^{N}(v) \left( \tilde{b}(v) - \rho \right)
.\]
Then \[
	G'(v) = (N-1) F^{N-1}(v) f(v) (v-\rho)
.\]
Integrating from $\rho$ to $v$ yields \[
	F^N(v) \left( \tilde{b}(v) - \rho \right)
	= \int_{\rho}^{v} (N-1)F^{N-1}(x) f(x) (x-\rho)dx
.\]
Since $\tilde{b}(\rho)=\rho$, this simplifies to \[
	\tilde{b}(v) - \rho = \frac{N-1}{N} m(v)
.\]
where \[
	m(v) = \E\left[X_{(1)} - \rho \mid X_{(1)} < v\right]
.\]

We can verify this $\tilde{b}$ is indeed increasing in $v$ and constitutes an equilibrium (by examining the sign of $\frac{du}{dr}$).
